% 87-02 ⑤(87쪽) — 보기 그래프: $x=-1$ 과 $x=3$ 에서 극소, 그 사이에서 극대인 아래로 볼록한 사차함수 개형. 그림 자체가 보기라 TikZ 로 다시 그림.
% 라벨은 원문 그대로 -1, O, 3, x, y, y=f(x) 와 점선 두 줄만.
\begin{tikzpicture}[baseline=(current bounding box.center), x=0.36cm, y=0.36cm, line width=0.5pt]
  \path (-2.9,-2.4) rectangle (5.15,3.5);
  \draw[->, >=stealth, line width=0.7pt] (-2.6,0) -- (4.85,0) node[below=1pt, font=\footnotesize] {$x$};
  \draw[->, >=stealth, line width=0.7pt] (0,-2.3) -- (0,3.0) node[left=1pt, font=\footnotesize] {$y$};
  \draw[densely dotted, line width=0.4pt] (-1,0) -- (-1,-1.3);
  \draw[densely dotted, line width=0.4pt] (3,0) -- (3,-1.65);
  \draw[line width=0.6pt, smooth] plot coordinates {(-2.15,2.6) (-2.05,1.9) (-1.95,1.3) (-1.85,0.7) (-1.75,0.25) (-1.6,-0.35) (-1.45,-0.8) (-1.25,-1.15) (-1,-1.3) (-0.8,-1.25) (-0.55,-1.1) (-0.3,-0.9) (0,-0.68) (0.3,-0.53) (0.6,-0.46) (0.85,-0.45) (1.1,-0.5) (1.4,-0.65) (1.8,-0.95) (2.2,-1.3) (2.6,-1.55) (3,-1.65) (3.3,-1.55) (3.6,-1.2) (3.85,-0.7) (4.05,-0.1) (4.2,0.6) (4.35,1.5) (4.45,2.3)};
  \node[above=1pt, font=\footnotesize] at (-1.1,0) {$-1$};
  \node[below left=-2pt and -3pt, font=\footnotesize] at (0,0) {$\pt{O}$};
  \node[above=1pt, font=\footnotesize] at (3,0) {$3$};
  \node[anchor=west, font=\footnotesize] at (0.25,2.75) {$y=f(x)$};
\end{tikzpicture}
